%%%%%%%%%%%%%%%%%%%%%%%%%%%%%%%%%%%%%%%%%%%%%%%%%%%%%%%%%%%%%%%%%%%%%%%%%%%%%%%%%%%
%%%  Filename: thesis_draft_bab1_bab2.tex
%%%  ---
%%%  Thesis Draft for NovaTrix RAG-Based DSS for ISO 27001 SoA Analysis
%%%  Compatible with thesisdtetiugm.cls
%%%  ---
%%%  Generated for: Muhammad Hilmi Dzaki Wismadi (22/497591/TK/54539)
%%%%%%%%%%%%%%%%%%%%%%%%%%%%%%%%%%%%%%%%%%%%%%%%%%%%%%%%%%%%%%%%%%%%%%%%%%%%%%%%%%%

%% Use option "bahasa" for Indonesian language
\documentclass[bachelor,bahasa]{thesisdtetiugm}

%======================================
% Information Input
%======================================
\author{Muhammad Hilmi Dzaki Wismadi}{22/497591/TK/54539}
\title{Perancangan dan Evaluasi Sistem Pendukung Keputusan Berbasis Retrieval-Augmented Generation untuk Analisis Statement of Applicability ISO/IEC 27001:2022 Annex A.5--A.8}
\program{Teknologi Informasi}{<<Program Coordinator>>}{<<NIP>>}
\departmenthead{<<Head of the Department>>}{<<NIP>>}
\major{Teknologi Informasi}
\yearsubmit{2026}
\examdate{<<Exam Date>>}

\addsupervisor{<<Dosen Pembimbing 1, S.T., M.Eng., Ph.D.>>}{<<NIP>>}
\addsupervisor{<<Dosen Pembimbing 2, S.T., M.Eng., Ph.D.>>}{<<NIP>>}

%======================================
% Additional Packages (if needed)
%======================================
\usepackage{longtable}
\usepackage{booktabs}
\usepackage{multirow}
\usepackage{array}
\usepackage{xcolor}

\begin{document}

%======================================
% COVER PAGE
%======================================
\printcover{sample/logougm.png}{Bachelor}

%======================================
% SDGs Declaration (Custom Section - to be placed appropriately)
%======================================
% Sustainable Development Goals:
% - SDG 9: Industry, Innovation and Infrastructure
% - SDG 16: Peace, Justice and Strong Institutions
% - SDG 4: Quality Education

%======================================
% TABLE OF CONTENTS, LIST OF TABLES, LIST OF FIGURES
%======================================
\cleardoublepage \phantomsection
\thetoc
\onehalfspacing
\tableofcontents
\singlespacing

\cleardoublepage \phantomsection
\thelot
\listoftables

\cleardoublepage \phantomsection
\thelof
\listoffigures

%======================================
% NOMENCLATURE / DAFTAR SINGKATAN
%======================================
\cleardoublepage \phantomsection
\chapter*{Daftar Singkatan}
\addcontentsline{toc}{chapter}{Daftar Singkatan}

\begin{tabular}{llp{3.5in}}
	AI & \hspace{1.5cm} = & \textit{Artificial Intelligence} (Kecerdasan Buatan) \\
	APAC & \hspace{1.5cm} = & \textit{Asia-Pacific} (Kawasan Asia-Pasifik) \\
	BSSN & \hspace{1.5cm} = & Badan Siber dan Sandi Negara \\
	CPU & \hspace{1.5cm} = & \textit{Central Processing Unit} \\
	DPR & \hspace{1.5cm} = & \textit{Dense Passage Retrieval} \\
	DSS & \hspace{1.5cm} = & \textit{Decision Support System} (Sistem Pendukung Keputusan) \\
	DTI & \hspace{1.5cm} = & Departemen Teknologi Informasi \\
	FAISS & \hspace{1.5cm} = & \textit{Facebook AI Similarity Search} \\
	ISMS & \hspace{1.5cm} = & \textit{Information Security Management System} \\
	ISO & \hspace{1.5cm} = & \textit{International Organization for Standardization} \\
	IEC & \hspace{1.5cm} = & \textit{International Electrotechnical Commission} \\
	JSON & \hspace{1.5cm} = & \textit{JavaScript Object Notation} \\
	LLM & \hspace{1.5cm} = & \textit{Large Language Model} \\
	NLP & \hspace{1.5cm} = & \textit{Natural Language Processing} \\
	RAG & \hspace{1.5cm} = & \textit{Retrieval-Augmented Generation} \\
	REST & \hspace{1.5cm} = & \textit{Representational State Transfer} \\
	RRF & \hspace{1.5cm} = & \textit{Reciprocal Rank Fusion} \\
	SoA & \hspace{1.5cm} = & \textit{Statement of Applicability} \\
	SPK & \hspace{1.5cm} = & Sistem Pendukung Keputusan \\
	UGM & \hspace{1.5cm} = & Universitas Gadjah Mada \\
\end{tabular}

%======================================
%  MAIN TEXT
%======================================
\startmain

%======================================
% BAB 1 - PENDAHULUAN
%======================================
\cleardoublepage \phantomsection
\chapter{Pendahuluan}
\label{chap:pendahuluan}

\section{Latar Belakang}
\label{sec:latar-belakang}

% Paragraf 1 - Urgensi Keamanan Siber
Perkembangan transformasi digital yang pesat dalam berbagai sektor telah meningkatkan ketergantungan organisasi terhadap sistem informasi dan infrastruktur teknologi. Namun demikian, peningkatan konektivitas dan digitalisasi tersebut juga membawa konsekuensi berupa meningkatnya risiko keamanan siber yang dihadapi oleh berbagai institusi. Laporan \textit{Check Point Global Cyber Attack Report} 2026 mencatat bahwa sektor pendidikan merupakan sektor yang paling banyak menjadi target serangan siber dengan rata-rata 4.749 serangan per organisasi per minggu, meningkat sebesar 7\% dibandingkan tahun sebelumnya \cite{checkpoint2026}. Secara regional, kawasan Asia-Pasifik (APAC) mencatat rata-rata 3.040 serangan siber per organisasi per minggu dengan peningkatan 3\% year-over-year \cite{checkpoint2026}. Data ini mengindikasikan bahwa institusi pendidikan, termasuk perguruan tinggi, menghadapi ancaman keamanan siber yang signifikan dan memerlukan pengelolaan keamanan informasi yang sistematis dan terstruktur.

% Paragraf 2 - Pentingnya ISO 27001
Dalam menghadapi tantangan keamanan siber tersebut, penerapan standar manajemen keamanan informasi menjadi kebutuhan yang mendesak bagi organisasi. ISO/IEC 27001 merupakan standar internasional yang menyediakan kerangka kerja \textit{Information Security Management System} (ISMS) berbasis risiko untuk membantu organisasi dalam melindungi aset informasinya \cite{iso27001review2022}. Standar ini memuat persyaratan sistematis untuk menetapkan, mengimplementasikan, memelihara, dan terus-menerus meningkatkan ISMS. Pada tanggal 25 Oktober 2022, ISO/IEC 27001:2022 diterbitkan sebagai pembaruan dari versi 2013, dengan pengurangan jumlah kontrol pada Annex A dari 114 menjadi 93 kontrol yang dikelompokkan ke dalam empat kategori: organisasional, manusia, fisik, dan teknologi \cite{iso27001review2022}. Pertumbuhan adopsi sertifikasi ISO 27001 secara global menunjukkan peningkatan signifikan dari 48.671 sertifikat valid pada 2018 menjadi 96.709 pada 2022, dengan lebih dari 40.000 organisasi tersertifikasi secara aktif \cite{frameworkcomparison2024}. Hal ini mengindikasikan bahwa ISO 27001 telah menjadi standar \textit{de facto} dalam pengelolaan keamanan informasi organisasi.

% Paragraf 3 - Keterbatasan Audit Manual (CRITICAL - Focus on audit process)
Meskipun ISO 27001 menyediakan kerangka kerja yang komprehensif, proses sertifikasi dan audit kepatuhan terhadap standar ini memerlukan sumber daya yang substansial. Berdasarkan studi empiris terhadap proses sertifikasi ISO 27001 di Portugal, durasi sertifikasi rata-rata berkisar antara 6 hingga 12 bulan dengan total biaya lebih dari EUR 50.000 untuk sebagian besar kasus \cite{automatingiso27001}. Rincian biaya tersebut meliputi jasa konsultan sebesar EUR 10.000--15.000, audit internal sebesar EUR 3.000--4.000, dan biaya akreditasi sebesar EUR 4.000--7.000 \cite{automatingiso27001}. Proses audit ISO 27001 terdiri dari beberapa tahapan yang memerlukan waktu intensif: fase perencanaan (\textit{planning}) yang berlangsung selama 2--4 minggu untuk menentukan ruang lingkup dan metodologi audit; fase implementasi (\textit{implementation}) yang memerlukan waktu 3--6 bulan untuk menerapkan kontrol keamanan dan mempersiapkan dokumentasi; serta fase audit (\textit{audit}) yang berlangsung selama 2--3 minggu untuk verifikasi kepatuhan oleh auditor eksternal \cite{automatingiso27001}. Kompleksitas proses ini diperparah oleh karakteristik audit manual yang bersifat \textit{labor-intensive}: pemetaan kontrol (\textit{control mapping}) terhadap 93 kontrol Annex A memerlukan analisis dokumen kebijakan, prosedur, dan bukti implementasi secara satu per satu; pengelolaan dokumentasi yang masih mengandalkan \textit{spreadsheet} dan dokumen teks menyebabkan inkonsistensi dan redundansi; serta beban kerja auditor yang tinggi dalam mengevaluasi \textit{Statement of Applicability} (SoA) yang memuat justifikasi untuk setiap kontrol \cite{csam2024, csam2validation2024}. Studi pada institusi pendidikan tinggi Kanada mencatat bahwa audit keamanan siber komprehensif memerlukan evaluasi terhadap 18 domain, 26 sub-domain, 87 \textit{checklist}, 169 kontrol, dan 429 sub-kontrol \cite{csam2024}.

% Paragraf 4 - LLM dan RAG dalam Audit (dengan tabel perbandingan)
Perkembangan teknologi \textit{Large Language Model} (LLM) dan arsitektur \textit{Retrieval-Augmented Generation} (RAG) membuka peluang untuk mengotomatisasi dan meningkatkan efisiensi proses audit kepatuhan. Beberapa penelitian terkini telah mengeksplorasi penerapan LLM dan RAG dalam domain audit dan kepatuhan regulasi. Tabel \ref{tab:llm-comparison} menyajikan perbandingan empat penelitian terkini dalam bidang ini.

\begin{table}[htbp]
\centering
\caption{Perbandingan Penelitian Sistem Kepatuhan Berbasis LLM dan RAG}
\label{tab:llm-comparison}
\small
\begin{tabular}{|p{2cm}|p{3cm}|p{3cm}|p{3cm}|p{3cm}|}
\hline
\textbf{Aspek} & \textbf{Paper \#23 (Gupta et al.)} & \textbf{Paper \#16 (LAJBR)} & \textbf{Paper \#18 (AuditGPT)} & \textbf{Paper \#21 (ReguQuery)} \\
\hline
\textbf{Fokus Utama} & \textit{Compliance checking} GDPR/CCPA dengan \textit{hybrid retrieval} & Pemetaan temuan audit terhadap dasar hukum pemerintahan & \textit{Pipeline} RAG untuk audit keuangan & \textit{Agentic framework} untuk kepatuhan regulasi keuangan India \\
\hline
\textbf{Metode \textit{Retrieval}} & Dua tahap: DPR (768-dim) + \textit{Cross-Encoder reranking} & Tiga jalur: BM25 \textit{sparse} + \textit{dense semantic} + \textit{weighted} RRF \textit{fusion} & \textit{Single-path dense retrieval} dengan LangChain & \textit{Single-path dense} (FAISS + Cohere \textit{embeddings}) \\
\hline
\textbf{Input--Output} & Dokumen kebijakan $\rightarrow$ Status kepatuhan + referensi GDPR & Teks temuan audit $\rightarrow$ Ketentuan hukum ter-\textit{ranking} (\textit{top-k}) & \textit{Query} + dokumen $\rightarrow$ \textit{Verdict} + bukti + rekomendasi & Dokumen perusahaan $\rightarrow$ JSON dengan skor (0--100), status, rekomendasi \\
\hline
\textbf{Kekuatan} & Presisi 98,2\% dengan \textit{false positive} rendah (6\%); LLM terkuantisasi untuk \textit{local deployment} & Ablasi rigor tinggi; \textit{multi-path fusion} meningkatkan Recall@5; metrik terukur (Recall@k, MRR@k) & \textit{Pipeline} lengkap dengan \textit{traceability} bukti; format \textit{verdict} terstruktur & Skor kepatuhan kuantitatif (0--100); deteksi \textit{irrelevance} untuk validasi cakupan \\
\hline
\textbf{Keterbatasan} & \textit{Cross-Encoder} menambah latensi (3,4s vs 0,75s DPR); fokus GDPR saja & Fokus komponen \textit{retrieval} saja, bukan sistem \textit{end-to-end}; tanpa luaran kepatuhan terstruktur & Evaluasi kualitatif; tidak ada ablasi \textit{retrieval} detail & Evaluasi kualitatif; \textit{cloud-based} (Cohere); tidak \textit{CPU-feasible} \\
\hline
\end{tabular}
\end{table}

Berdasarkan tinjauan literatur pada Tabel \ref{tab:llm-comparison}, Paper \#23 (Gupta et al., 2025) mengembangkan sistem \textit{compliance checking} untuk regulasi GDPR dan CCPA dengan menggabungkan \textit{Dense Passage Retrieval} (DPR) dan model \textit{Cross-Encoder} dalam arsitektur \textit{hybrid retrieval}, mencapai presisi 98,2\% meskipun dengan peningkatan latensi \cite{gupta2025compliance}. Berbeda dengan pendekatan tersebut, Paper \#16 (LAJBR) mengadopsi pendekatan \textit{multi-path retrieval} untuk pemetaan temuan audit terhadap dasar hukum pada domain pemerintahan, mengintegrasikan tiga jalur \textit{retrieval} yang difusikan menggunakan \textit{weighted Reciprocal Rank Fusion} (RRF) \cite{lajbr2025}. Sementara itu, Paper \#18 (AuditGPT) menyajikan \textit{pipeline} RAG yang lebih lengkap untuk domain audit keuangan dengan luaran berupa \textit{verdict} kepatuhan yang disertai kutipan bukti yang dapat ditelusuri ke dokumen sumber \cite{auditgpt2025}. Pendekatan yang lebih komprehensif dari sisi sistem ditunjukkan oleh Paper \#21 (ReguQuery), yang mengembangkan \textit{agentic compliance framework} dengan luaran terstruktur berupa skor kepatuhan kuantitatif (0--100) dalam format JSON \cite{reguquery2025}. Dari keempat paper tersebut, dapat diidentifikasi bahwa penelitian yang ada cenderung berfokus secara parsial: Paper \#23 dan \#16 unggul pada aspek \textit{retrieval} namun tidak menyediakan sistem audit \textit{end-to-end}; Paper \#18 menyajikan \textit{pipeline} lengkap namun kurang rigoros dalam ablasi komponen \textit{retrieval}; sedangkan Paper \#21 menawarkan luaran terstruktur namun bergantung pada infrastruktur \textit{cloud}.

% Paragraf 5 - Kebutuhan Sistem NovaTrix
Celah penelitian tersebut membuka peluang untuk pengembangan sistem yang mengintegrasikan strategi \textit{retrieval} yang optimal dengan \textit{pipeline} audit \textit{end-to-end} yang menghasilkan luaran terstruktur---dalam konteks standar keamanan informasi ISO/IEC 27001:2022---dengan mempertimbangkan keterbatasan sumber daya komputasi lokal. Proses analisis \textit{Statement of Applicability} (SoA) yang memerlukan pemetaan setiap temuan audit terhadap 93 kontrol Annex A beserta keputusan \textit{applicability} dan status implementasi menjadi kandidat yang tepat untuk otomatisasi berbasis RAG. Penelitian ini mengusulkan pengembangan NovaTrix, sebuah sistem pendukung keputusan (\textit{decision support system}) berbasis RAG untuk analisis SoA ISO/IEC 27001:2022 dengan fokus pada Annex A.5--A.8. Arsitektur sistem dirancang dengan mempertimbangkan keterbatasan perangkat keras CPU-\textit{only} (Intel i7-12550U, 16GB RAM, tanpa GPU) melalui penggunaan model LLM terkuantisasi yang dijalankan secara lokal menggunakan Ollama, \textit{embedding} model ringan (paraphrase-multilingual-MiniLM-L12-v2, 384-dim), dan \textit{vector store} FAISS dengan \textit{IndexFlatIP}. Sistem ini diimplementasikan sebagai aplikasi web berbasis React untuk \textit{frontend} dan Express untuk \textit{backend}, dengan luaran terstruktur berupa \texttt{control\_id}, \texttt{applicable} (Yes/No), \texttt{implementation\_status} (Implemented/Partially Implemented/Not Implemented), \texttt{justification}, dan \texttt{recommendation}.

\section{Rumusan Masalah}
\label{sec:rumusan-masalah}

Berdasarkan latar belakang yang telah diuraikan, rumusan masalah dalam penelitian ini adalah sebagai berikut:

\begin{enumerate}
	\item Bagaimana merancang arsitektur sistem pendukung keputusan berbasis \textit{Retrieval-Augmented Generation} (RAG) yang dapat menganalisis \textit{Statement of Applicability} (SoA) ISO/IEC 27001:2022 Annex A.5--A.8 dengan keterbatasan sumber daya komputasi lokal?
	
	\item Bagaimana kinerja sistem RAG dibandingkan dengan pendekatan LLM-\textit{only} (\textit{baseline}) dalam hal akurasi pemetaan kontrol, ketepatan keputusan \textit{applicability}, dan ketepatan penentuan status implementasi?
	
	\item Bagaimana mengevaluasi efektivitas komponen \textit{retrieval} sistem RAG melalui metrik Recall@k dan analisis kualitas konteks yang diambil?
\end{enumerate}

\section{Tujuan Penelitian}
\label{sec:tujuan-penelitian}

Berdasarkan rumusan masalah yang telah diidentifikasi, tujuan penelitian ini adalah sebagai berikut:

\begin{enumerate}
	\item Merancang dan mengimplementasikan arsitektur sistem pendukung keputusan berbasis RAG untuk analisis SoA ISO/IEC 27001:2022 Annex A.5--A.8 yang dapat dijalankan pada perangkat keras dengan keterbatasan sumber daya (CPU-\textit{only}, 16GB RAM).
	
	\item Mengevaluasi kinerja sistem RAG secara komparatif terhadap pendekatan LLM-\textit{only} menggunakan metrik akurasi pemetaan kontrol, akurasi keputusan \textit{applicability}, akurasi penentuan status implementasi, serta presisi, \textit{recall}, dan F1-\textit{score}.
	
	\item Menganalisis efektivitas komponen \textit{retrieval} melalui pengukuran Recall@k (k=1, 3, 5) dan evaluasi kualitas konteks yang diambil untuk mendukung pengambilan keputusan.
	
	\item Mengukur performa operasional sistem meliputi waktu respons (\textit{latency}) dan stabilitas luaran untuk memastikan kelayakan penggunaan praktis.
	
	\item Mengidentifikasi pola kesalahan dan keterbatasan sistem melalui analisis kasus-kasus salah klasifikasi untuk memberikan rekomendasi perbaikan.
\end{enumerate}

\section{Batasan Penelitian}
\label{sec:batasan-penelitian}

Penelitian ini memiliki batasan-batasan sebagai berikut:

\begin{enumerate}
	\item \textbf{Cakupan Kontrol ISO 27001:2022}: Penelitian ini membatasi analisis pada kontrol Annex A.5 (Organizational Controls), A.6 (People Controls), A.7 (Physical Controls), dan A.8 (Technological Controls), yang mencakup total 93 kontrol keamanan informasi.
	
	\item \textbf{Keterbatasan Perangkat Keras}: Sistem dirancang untuk berjalan pada lingkungan CPU-\textit{only} dengan spesifikasi Intel i7-12550U, 16GB RAM, dan sistem operasi Windows 11, tanpa akselerasi GPU. Hal ini membatasi penggunaan model LLM pada varian 7B--8B terkuantisasi.
	
	\item \textbf{Model LLM}: Penelitian menggunakan model qwen2.5:3b yang dijalankan secara lokal melalui Ollama dengan parameter \texttt{temperature=0} untuk determinisme luaran.
	
	\item \textbf{Model \textit{Embedding}}: Penelitian menggunakan model paraphrase-multilingual-MiniLM-L12-v2 (384 dimensi) yang mendukung teks berbahasa Indonesia.
	
	\item \textbf{\textit{Vector Store}}: Penelitian menggunakan FAISS dengan \textit{IndexFlatIP} (pencarian eksak berbasis \textit{inner product}) tanpa \textit{hybrid search} atau \textit{reranking}.
	
	\item \textbf{Dataset Evaluasi}: Evaluasi dilakukan menggunakan dataset \textit{gold standard} yang disusun secara manual, berisi minimal 100 kalimat temuan audit yang dilabeli dengan \texttt{control\_id}, \texttt{applicable}, \texttt{implementation\_status}, dan \texttt{justification}.
	
	\item \textbf{Bahasa}: Input sistem berupa teks temuan audit dalam bahasa Indonesia.
	
	\item \textbf{Konteks Studi Kasus}: Penelitian ini menggunakan konteks institusi pendidikan tinggi (Departemen Teknologi Informasi, Universitas Gadjah Mada) sebagai domain aplikasi.
\end{enumerate}

\section{Manfaat Penelitian}
\label{sec:manfaat-penelitian}

Penelitian ini diharapkan memberikan manfaat baik secara akademis maupun praktis sebagai berikut:

\subsection{Manfaat Akademis}

\begin{enumerate}
	\item Memberikan kontribusi empiris berupa perbandingan kuantitatif antara pendekatan RAG dan LLM-\textit{only} untuk tugas pemetaan kontrol keamanan informasi ISO 27001, yang dapat menjadi referensi bagi penelitian selanjutnya di bidang otomatisasi audit berbasis \textit{natural language processing}.
	
	\item Menyediakan kerangka eksperimental yang terstruktur dan dapat direproduksi untuk evaluasi sistem RAG pada domain kepatuhan regulasi, mencakup metrik klasifikasi (akurasi, presisi, \textit{recall}, F1-\textit{score}), metrik \textit{retrieval} (Recall@k), dan metrik operasional (latensi).
	
	\item Mengidentifikasi celah penelitian dan memberikan rekomendasi untuk pengembangan sistem RAG berbasis LLM dalam konteks audit keamanan informasi, khususnya terkait optimasi \textit{retrieval} dan penanganan keterbatasan sumber daya komputasi.
\end{enumerate}

\subsection{Manfaat Praktis}

\begin{enumerate}
	\item Menyediakan prototipe sistem pendukung keputusan berbasis RAG (NovaTrix) yang dapat digunakan oleh praktisi keamanan informasi dan auditor ISO 27001 untuk mempercepat proses analisis \textit{Statement of Applicability} dan pemetaan kontrol.
	
	\item Mendemonstrasikan kelayakan penerapan teknologi LLM dan RAG pada lingkungan dengan keterbatasan sumber daya (\textit{resource-constrained environment}) yang umum ditemui pada organisasi pendidikan dan usaha kecil menengah.
	
	\item Memberikan panduan teknis implementasi sistem RAG untuk domain kepatuhan regulasi yang dapat diadaptasi untuk standar keamanan informasi lainnya atau kerangka kerja kepatuhan serupa.
\end{enumerate}

\section{Sistematika Penulisan}
\label{sec:sistematika-penulisan}

Sistematika penulisan tugas akhir ini terdiri dari enam bab yang disusun sebagai berikut:

\textbf{Bab I Pendahuluan} membahas latar belakang penelitian yang mencakup urgensi keamanan siber, pentingnya ISO 27001, keterbatasan audit manual, potensi LLM dan RAG dalam audit, serta kebutuhan pengembangan sistem NovaTrix. Bab ini juga menyajikan rumusan masalah, tujuan penelitian, batasan penelitian, dan manfaat penelitian.

\textbf{Bab II Tinjauan Pustaka dan Dasar Teori} menyajikan tinjauan pustaka berupa kajian terhadap penelitian-penelitian terdahulu yang relevan dengan sistem kepatuhan berbasis LLM dan RAG, serta dasar teori yang mencakup konsep ISO/IEC 27001, ISMS, \textit{Statement of Applicability}, \textit{Retrieval-Augmented Generation}, \textit{Large Language Model}, audit keamanan informasi, dan sistem pendukung keputusan.

\textbf{Bab III Metodologi Penelitian} menjelaskan metodologi yang digunakan dalam penelitian, meliputi desain eksperimen komparatif, arsitektur sistem, \textit{pipeline} RAG, prosedur pengumpulan data, dan metrik evaluasi.

\textbf{Bab IV Implementasi Sistem} membahas implementasi teknis sistem NovaTrix, mencakup struktur proyek, komponen \textit{embedding} dan \textit{indexing}, modul \textit{retrieval}, integrasi LLM, dan antarmuka pengguna.

\textbf{Bab V Hasil dan Pembahasan} menyajikan hasil eksperimen komparatif antara sistem RAG dan \textit{baseline} LLM-\textit{only}, analisis metrik klasifikasi dan \textit{retrieval}, analisis kasus kesalahan, serta pembahasan implikasi hasil penelitian.

\textbf{Bab VI Kesimpulan dan Saran} merangkum kesimpulan dari hasil penelitian, keterbatasan yang ditemui, serta saran untuk penelitian dan pengembangan di masa mendatang.

%======================================
% BAB 2 - TINJAUAN PUSTAKA DAN DASAR TEORI
%======================================
\cleardoublepage \phantomsection
\chapter{Tinjauan Pustaka dan Dasar Teori}
\label{chap:tinjauan-pustaka}

\section{Tinjauan Pustaka}
\label{sec:tinjauan-pustaka}

Bab ini menyajikan tinjauan terhadap penelitian-penelitian terdahulu yang relevan dengan pengembangan sistem kepatuhan dan audit berbasis \textit{Large Language Model} (LLM) dan arsitektur \textit{Retrieval-Augmented Generation} (RAG). Tinjauan pustaka ini bertujuan untuk mengidentifikasi celah penelitian yang menjadi dasar pengembangan sistem NovaTrix.

\subsection{Sistem Kepatuhan Berbasis LLM dan RAG}

Pemanfaatan LLM dan RAG untuk otomatisasi pemeriksaan kepatuhan regulasi telah menarik perhatian peneliti dalam beberapa tahun terakhir. Gupta et al. (2025) mengembangkan sistem \textit{compliance checking} untuk regulasi GDPR dan CCPA dengan menggabungkan \textit{Dense Passage Retrieval} (DPR) dan model \textit{Cross-Encoder} dalam arsitektur \textit{hybrid retrieval} \cite{gupta2025compliance}. Pendekatan dua tahap ini---DPR untuk \textit{initial retrieval} diikuti \textit{Cross-Encoder} untuk \textit{reranking}---terbukti mencapai presisi 98,2\% dengan tingkat \textit{false positive} yang rendah (6\%), meskipun dengan konsekuensi peningkatan latensi dari 0,75 detik menjadi 3,4 detik. Penggunaan LLaMA 3.1 terkuantisasi (8-bit) pada penelitian ini menunjukkan kelayakan \textit{deployment} lokal untuk sistem kepatuhan berbasis LLM.

Berbeda dengan pendekatan Gupta et al. yang berfokus pada optimasi \textit{reranking}, penelitian LAJBR mengadopsi pendekatan \textit{multi-path retrieval} untuk pemetaan temuan audit terhadap dasar hukum pada domain pemerintahan \cite{lajbr2025}. Sistem ini mengintegrasikan tiga jalur \textit{retrieval}---BM25 \textit{sparse}, \textit{dense semantic} dengan klausul hipotetikal, dan \textit{dense semantic} dengan \textit{legal-style rewrite}---yang kemudian difusikan menggunakan \textit{weighted Reciprocal Rank Fusion} (RRF). Studi ablasi yang dilakukan menunjukkan bahwa kombinasi ketiga jalur secara konsisten meningkatkan Recall@5 dan MRR@5 dibandingkan penggunaan jalur tunggal. Namun demikian, penelitian ini membatasi fokusnya pada komponen \textit{retrieval} tanpa menyediakan luaran kepatuhan terstruktur yang dapat langsung digunakan dalam alur kerja audit.

Pendekatan yang lebih lengkap dari sisi \textit{pipeline} end-to-end disajikan oleh AuditGPT untuk domain audit keuangan \cite{auditgpt2025}. Sistem ini tidak hanya melakukan \textit{retrieval} konteks yang relevan, tetapi juga menghasilkan \textit{verdict} kepatuhan (\textit{Compliant}, \textit{Non-Compliant}, atau \textit{Requires Review}) yang disertai dengan kutipan bukti (\textit{evidence excerpts}) yang dapat ditelusuri ke dokumen sumber. Integrasi dengan LangChain sebagai \textit{orchestrator} memungkinkan alur pemrosesan yang terstruktur dari \textit{ingestion} hingga generasi rekomendasi. Meskipun demikian, evaluasi yang disajikan cenderung kualitatif dan tidak menyediakan ablasi detail terhadap komponen \textit{retrieval}.

ReguQuery mengambil pendekatan yang lebih komprehensif dari sisi sistem dengan mengembangkan \textit{agentic compliance framework} untuk sektor keuangan India (SEBI, RBI, AMFI) \cite{reguquery2025}. Keunggulan utama ReguQuery terletak pada luaran terstruktur berupa skor kepatuhan kuantitatif (0-100), daftar seksi yang tidak patuh, serta rekomendasi perbaikan dalam format JSON. Framework ini juga mengimplementasikan \textit{irrelevance detection} untuk memvalidasi apakah dokumen input sesuai dengan cakupan regulasi yang dianalisis. Namun, ketergantungan pada layanan \textit{cloud-based} (Cohere Command-R) dan evaluasi yang sebagian besar kualitatif menjadi keterbatasan untuk adopsi pada lingkungan dengan keterbatasan sumber daya komputasi.

\subsection{Penerapan LLM dalam Audit ISO 27001}

Dalam konteks spesifik ISO 27001, beberapa penelitian telah mengeksplorasi potensi LLM untuk mendukung proses audit dan kepatuhan. Studi kelayakan oleh Rathbun et al. (2024) mengevaluasi tiga model LLM---Claude Sonnet 3.5, ChatGPT 4o, dan Llama 3.1---untuk tiga tugas audit ISO 27001: generasi SoA, perencanaan ruang lingkup audit, dan pengembangan \textit{checklist} \cite{llmiso27001pilot}. Evaluasi menggunakan skala Likert 0--3 menunjukkan bahwa Claude Sonnet 3.5 mencapai skor tertinggi untuk relevansi dan inferensi, meskipun skor \textit{completeness} pada set pertama masih rendah. Temuan ini mengindikasikan bahwa LLM memiliki potensi untuk mempercepat alur kerja kepatuhan, namun memerlukan penanganan khusus untuk memastikan kelengkapan dan akurasi luaran.

Penelitian lain oleh Marques et al. (2024) melakukan studi kelayakan otomatisasi konsultansi dan audit ISO 27001 menggunakan AI melalui wawancara dengan enam pakar sertifikasi \cite{automatingiso27001}. Studi ini mengidentifikasi bahwa proses sertifikasi rata-rata memerlukan waktu 6--12 bulan dengan biaya lebih dari EUR 50.000, mengindikasikan kebutuhan mendesak akan otomatisasi. Pakar yang diwawancarai mengakui potensi AI untuk mempercepat analisis dokumen dan pemetaan kontrol, namun menekankan bahwa validasi manusia tetap diperlukan untuk memastikan akurasi dan akuntabilitas.

\subsection{Model Maturitas dan Audit Keamanan Siber di Institusi Pendidikan}

Penelitian tentang audit keamanan siber di institusi pendidikan tinggi memberikan konteks penting untuk pengembangan sistem NovaTrix. Yigit Ozkan dan Spruit (2024) mengembangkan CyberSecurity Audit Model (CSAM) yang terdiri dari 18 domain, 26 sub-domain, 87 \textit{checklist}, 169 kontrol, dan 429 sub-kontrol \cite{csam2024}. Validasi CSAM pada institusi pendidikan tinggi Kanada dengan lebih dari 2.500 karyawan dan 15.000 mahasiswa menunjukkan tingkat maturitas keamanan siber sebesar 80\% (kategori \textit{Mature}, band 71--90\%). Temuan ini menegaskan kompleksitas audit keamanan siber yang komprehensif dan kebutuhan akan alat bantu yang dapat menyederhanakan proses evaluasi.

Validasi lanjutan CSAM 2.0 oleh Yigit Ozkan et al. (2025) mencakup tiga skenario: audit domain tunggal, audit multi-domain terpilih, dan audit seluruh domain \cite{csam2validation2024}. Hasil menunjukkan variasi skor yang signifikan antar domain, dengan domain Legal/Compliance mencapai 100\%, sementara domain IoT hanya mencapai 50\%. Temuan ini mengindikasikan bahwa sistem pendukung keputusan untuk audit keamanan perlu mampu mengakomodasi variabilitas tingkat kepatuhan antar domain kontrol.

\subsection{Keterbatasan LLM dan Solusi RAG}

Meskipun LLM menunjukkan potensi besar dalam berbagai tugas NLP, model ini memiliki keterbatasan inheren yang perlu diatasi untuk aplikasi audit dan kepatuhan. Survei oleh Ji et al. (2023) mengidentifikasi bahwa LLM prone terhadap halusinasi, yaitu generasi konten yang menyimpang dari fakta atau berisi informasi fabrikasi \cite{hallucinationsurvey}. Evaluasi FACTSCORE menunjukkan bahwa ChatGPT hanya mencapai akurasi faktual 58\% pada tugas generasi biografi \cite{factscore2023}. HaluEval benchmark menemukan bahwa ChatGPT menghasilkan konten berhalusinasi pada sekitar 19,5\% respons \cite{halueval2023}. Keterbatasan ini menjadi kritis dalam konteks audit kepatuhan yang memerlukan akurasi tinggi.

\textit{Retrieval-Augmented Generation} (RAG) telah muncul sebagai solusi yang menjanjikan untuk mengatasi keterbatasan LLM. Lewis et al. (2020) memperkenalkan arsitektur RAG yang menggabungkan memori parametrik (LLM) dan non-parametrik (retrieval) untuk generasi bahasa \cite{ragoriginal}. RAG terbukti menghasilkan luaran yang lebih spesifik, beragam, dan faktual dibandingkan baseline seq2seq parametrik. Survei komprehensif oleh Gao et al. (2024) mengkategorikan evolusi paradigma RAG dari \textit{Naive RAG} ke \textit{Advanced RAG} hingga \textit{Modular RAG}, yang memungkinkan pembaruan pengetahuan berkelanjutan dan kustomisasi domain-spesifik tanpa \textit{retraining} \cite{ragsurvey2024}.

\subsection{Celah Penelitian dan Posisi Penelitian Ini}

Berdasarkan tinjauan pustaka yang telah dilakukan, dapat diidentifikasi beberapa celah penelitian:

\begin{enumerate}
	\item Penelitian sistem kepatuhan berbasis RAG yang ada cenderung berfokus pada domain GDPR/CCPA, audit keuangan, atau regulasi sektor keuangan, belum ada yang secara spesifik menangani pemetaan kontrol ISO/IEC 27001:2022 Annex A.
	
	\item Sistem yang ada umumnya memerlukan infrastruktur \textit{cloud} atau GPU untuk inferensi LLM, sehingga tidak dapat diadopsi oleh organisasi dengan keterbatasan sumber daya komputasi.
	
	\item Evaluasi kuantitatif yang membandingkan pendekatan RAG vs LLM-\textit{only} secara terkontrol (dengan variabel konstan yang jelas) masih terbatas, khususnya untuk domain audit keamanan informasi.
	
	\item Luaran sistem audit yang ada belum seluruhnya terstruktur untuk mendukung pembuatan SoA secara langsung.
\end{enumerate}

Penelitian ini mengisi celah tersebut dengan mengembangkan NovaTrix, sistem pendukung keputusan berbasis RAG untuk analisis SoA ISO/IEC 27001:2022 yang: (1) berfokus pada 93 kontrol Annex A.5--A.8; (2) dirancang untuk lingkungan CPU-\textit{only}; (3) dievaluasi secara kuantitatif dengan eksperimen komparatif terkontrol; dan (4) menghasilkan luaran terstruktur yang langsung mendukung pembuatan SoA.

\section{Dasar Teori}
\label{sec:dasar-teori}

Bagian ini menyajikan landasan teori yang mendasari pengembangan sistem NovaTrix, mencakup konsep-konsep fundamental terkait keamanan informasi, standar ISO/IEC 27001, teknologi \textit{Large Language Model}, dan arsitektur \textit{Retrieval-Augmented Generation}.

\subsection{ISO/IEC 27001:2022}

ISO/IEC 27001 merupakan standar internasional yang menetapkan persyaratan untuk sistem manajemen keamanan informasi (\textit{Information Security Management System}/ISMS) dalam konteks organisasi \cite{iso27001review2022}. Standar ini diterbitkan bersama oleh \textit{International Organization for Standardization} (ISO) dan \textit{International Electrotechnical Commission} (IEC). Versi terbaru, ISO/IEC 27001:2022, diterbitkan pada 25 Oktober 2022 sebagai pembaruan dari versi 2013.

ISO/IEC 27001:2022 menyediakan kerangka kerja berbasis risiko untuk melindungi kerahasiaan (\textit{confidentiality}), integritas (\textit{integrity}), dan ketersediaan (\textit{availability}) informasi organisasi. Standar ini terdiri dari dua bagian utama: klausul utama (4--10) yang menetapkan persyaratan ISMS, dan Annex A yang memuat katalog kontrol keamanan informasi. Annex A versi 2022 memuat 93 kontrol yang dikelompokkan ke dalam empat kategori:

\begin{itemize}
	\item \textbf{A.5 Organizational Controls} (37 kontrol): Kontrol yang terkait dengan kebijakan, prosedur, dan tata kelola organisasi.
	\item \textbf{A.6 People Controls} (8 kontrol): Kontrol yang terkait dengan sumber daya manusia, termasuk \textit{screening}, kesadaran keamanan, dan tanggung jawab.
	\item \textbf{A.7 Physical Controls} (14 kontrol): Kontrol yang terkait dengan keamanan fisik, termasuk perimeter keamanan, akses fisik, dan perlindungan peralatan.
	\item \textbf{A.8 Technological Controls} (34 kontrol): Kontrol yang terkait dengan teknologi informasi, termasuk kontrol akses, kriptografi, dan keamanan jaringan.
\end{itemize}

\subsection{Information Security Management System (ISMS)}

\textit{Information Security Management System} (ISMS) adalah pendekatan sistematis untuk mengelola informasi sensitif organisasi agar tetap aman \cite{iso27001review2022}. ISMS mencakup orang, proses, dan sistem TI dengan menerapkan proses manajemen risiko. Komponen utama ISMS meliputi:

\begin{enumerate}
	\item \textbf{Kebijakan Keamanan Informasi}: Dokumen tingkat tinggi yang menetapkan komitmen organisasi terhadap keamanan informasi.
	\item \textbf{Penilaian Risiko}: Proses identifikasi, analisis, dan evaluasi risiko keamanan informasi.
	\item \textbf{Penanganan Risiko}: Pemilihan dan implementasi kontrol untuk mengurangi risiko ke tingkat yang dapat diterima.
	\item \textbf{Statement of Applicability (SoA)}: Dokumen yang mencantumkan semua kontrol dari Annex A beserta keputusan \textit{applicability} dan justifikasinya.
	\item \textbf{Pemantauan dan Tinjauan}: Proses berkelanjutan untuk memastikan efektivitas ISMS.
\end{enumerate}

\subsection{Statement of Applicability (SoA)}

\textit{Statement of Applicability} (SoA) merupakan dokumen wajib dalam ISO/IEC 27001 yang berfungsi sebagai penghubung antara penilaian risiko dan implementasi kontrol keamanan \cite{iso27001review2022}. SoA memuat daftar semua kontrol dari Annex A beserta:

\begin{itemize}
	\item \textbf{Keputusan \textit{Applicability}}: Apakah kontrol tersebut berlaku (Yes) atau tidak berlaku (No) untuk organisasi.
	\item \textbf{Justifikasi}: Alasan mengapa kontrol diterapkan atau dikecualikan.
	\item \textbf{Status Implementasi}: Tingkat implementasi kontrol (Implemented, Partially Implemented, Not Implemented, atau Planned).
	\item \textbf{Referensi Bukti}: Dokumen, kebijakan, atau prosedur yang mendukung implementasi kontrol.
\end{itemize}

Penyusunan SoA memerlukan analisis mendalam terhadap setiap kontrol dalam konteks organisasi, yang menjadikannya kandidat ideal untuk otomatisasi berbasis sistem cerdas.

\subsection{Large Language Model (LLM)}

\textit{Large Language Model} (LLM) adalah model bahasa berbasis arsitektur \textit{transformer} yang dilatih pada korpus teks skala besar untuk mempelajari representasi statistik bahasa \cite{llmsurvey2023}. Karakteristik utama LLM meliputi:

\begin{enumerate}
	\item \textbf{Skala Parameter}: Model modern memiliki miliaran hingga ratusan miliar parameter, memungkinkan kapasitas pengetahuan yang besar.
	\item \textbf{Kemampuan Emergent}: Kemampuan seperti \textit{chain-of-thought reasoning} dan \textit{in-context learning} muncul pada model dengan parameter lebih dari 100 miliar.
	\item \textbf{Generalisasi Zero-Shot}: Kemampuan untuk menyelesaikan tugas baru tanpa \textit{fine-tuning} khusus.
	\item \textbf{Kuantisasi}: Teknik kompresi model yang memungkinkan inferensi pada perangkat keras terbatas dengan trade-off akurasi minimal.
\end{enumerate}

Dalam konteks penelitian ini, digunakan model qwen2.5:3b yang dijalankan secara lokal melalui Ollama, memanfaatkan kuantisasi untuk memungkinkan inferensi pada lingkungan CPU-\textit{only}.

\subsection{Retrieval-Augmented Generation (RAG)}

\textit{Retrieval-Augmented Generation} (RAG) merupakan arsitektur yang menggabungkan komponen \textit{retrieval} (pengambilan informasi) dengan komponen \textit{generation} (generasi teks) untuk menghasilkan luaran yang lebih akurat dan terkini \cite{ragoriginal, ragsurvey2024}. Komponen utama arsitektur RAG meliputi:

\begin{enumerate}
	\item \textbf{Knowledge Base}: Koleksi dokumen atau informasi yang menjadi sumber pengetahuan eksternal.
	\item \textbf{Embedding Model}: Model yang mengubah teks menjadi representasi vektor untuk pencarian semantik.
	\item \textbf{Vector Store}: Basis data vektor yang menyimpan \textit{embedding} dokumen dan mendukung pencarian \textit{nearest neighbor}.
	\item \textbf{Retriever}: Komponen yang mengambil dokumen relevan dari \textit{vector store} berdasarkan query pengguna.
	\item \textbf{Generator}: Model bahasa yang menghasilkan respons berdasarkan query dan konteks yang diambil.
\end{enumerate}

Keunggulan RAG dibandingkan LLM-\textit{only} meliputi:
\begin{itemize}
	\item Mengurangi halusinasi dengan menyediakan konteks faktual yang relevan.
	\item Memungkinkan pembaruan pengetahuan tanpa \textit{retraining} model.
	\item Meningkatkan transparansi melalui \textit{traceability} sumber informasi.
	\item Memungkinkan kustomisasi domain-spesifik melalui \textit{knowledge base} yang dapat disesuaikan.
\end{itemize}

\subsection{Audit Keamanan Informasi}

Audit keamanan informasi adalah proses sistematis untuk mengevaluasi efektivitas kontrol keamanan informasi organisasi terhadap standar atau kriteria yang ditetapkan \cite{csam2024}. Tipe audit keamanan informasi meliputi:

\begin{enumerate}
	\item \textbf{Audit Internal}: Dilakukan oleh auditor internal organisasi untuk mengevaluasi kepatuhan dan efektivitas ISMS.
	\item \textbf{Audit Eksternal (Sertifikasi)}: Dilakukan oleh badan sertifikasi independen untuk memverifikasi kepatuhan terhadap standar ISO 27001.
	\item \textbf{Audit Gap}: Evaluasi pendahuluan untuk mengidentifikasi kesenjangan antara kondisi saat ini dan persyaratan standar.
\end{enumerate}

Proses audit umumnya mencakup tahapan: perencanaan, pengumpulan bukti, evaluasi, dan pelaporan. Setiap tahapan memerlukan analisis dokumen, wawancara, dan observasi yang intensif.

\subsection{Sistem Pendukung Keputusan (Decision Support System)}

Sistem Pendukung Keputusan (\textit{Decision Support System}/DSS) adalah sistem informasi berbasis komputer yang mendukung aktivitas pengambilan keputusan bisnis atau organisasi \cite{dsstheory}. Komponen utama DSS meliputi:

\begin{enumerate}
	\item \textbf{Subsistem Data}: Basis data yang menyimpan informasi relevan untuk pengambilan keputusan.
	\item \textbf{Subsistem Model}: Model analitik atau algoritma yang memproses data untuk menghasilkan insight.
	\item \textbf{Subsistem Dialog}: Antarmuka pengguna yang memungkinkan interaksi dengan sistem.
	\item \textbf{Subsistem Pengetahuan}: Basis pengetahuan yang menyimpan aturan atau heuristik keputusan.
\end{enumerate}

Dalam konteks NovaTrix, DSS berbasis RAG mengintegrasikan:
\begin{itemize}
	\item \textit{Knowledge base} ISO 27001 sebagai subsistem pengetahuan.
	\item \textit{Pipeline} RAG sebagai subsistem model.
	\item Antarmuka web sebagai subsistem dialog.
	\item Basis data temuan audit dan kontrol sebagai subsistem data.
\end{itemize}

\subsection{Vector Database dan Semantic Search}

\textit{Vector database} adalah basis data yang dioptimalkan untuk menyimpan dan melakukan pencarian pada data vektor berdimensi tinggi \cite{vectordbsurvey}. Dalam konteks RAG, \textit{vector database} digunakan untuk menyimpan \textit{embedding} dokumen dan melakukan \textit{semantic search} berdasarkan similaritas vektor.

FAISS (\textit{Facebook AI Similarity Search}) merupakan library yang digunakan dalam penelitian ini, mendukung berbagai tipe indeks:
\begin{itemize}
	\item \textbf{IndexFlatIP}: Pencarian eksak berbasis \textit{inner product}, memberikan akurasi tertinggi dengan kompleksitas O(n).
	\item \textbf{IndexFlatL2}: Pencarian eksak berbasis jarak Euclidean.
	\item \textbf{IndexIVF}: Pencarian aproksimasi dengan \textit{inverted file index} untuk dataset besar.
\end{itemize}

Penelitian ini menggunakan IndexFlatIP karena jumlah kontrol ISO 27001 yang relatif kecil (93 kontrol) tidak memerlukan indeks aproksimasi.

\subsection{Metrik Evaluasi Sistem Klasifikasi dan Retrieval}

Evaluasi sistem NovaTrix menggunakan dua kategori metrik:

\textbf{Metrik Klasifikasi:}
\begin{itemize}
	\item \textbf{Akurasi}: Proporsi prediksi yang benar dari total prediksi.
	\item \textbf{Presisi}: Proporsi prediksi positif yang benar dari total prediksi positif.
	\item \textbf{Recall}: Proporsi prediksi positif yang benar dari total aktual positif.
	\item \textbf{F1-Score}: Rata-rata harmonik presisi dan \textit{recall}.
\end{itemize}

\textbf{Metrik Retrieval:}
\begin{itemize}
	\item \textbf{Recall@k}: Proporsi query di mana dokumen relevan muncul dalam top-k hasil \textit{retrieval}.
	\item \textbf{MRR (Mean Reciprocal Rank)}: Rata-rata invers peringkat dokumen relevan pertama.
\end{itemize}

%======================================
% REFERENCES
%======================================
\cleardoublepage \phantomsection
\thereferences
\bibliography{references}

% Placeholder citations for BibTeX
% NOTE: The actual .bib file needs to be updated with these entries:
% @misc{checkpoint2026, ... }
% @article{iso27001review2022, ... }
% @article{automatingiso27001, ... }
% @inproceedings{gupta2025compliance, ... }
% @inproceedings{lajbr2025, ... }
% @inproceedings{auditgpt2025, ... }
% @inproceedings{reguquery2025, ... }
% @inproceedings{llmiso27001pilot, ... }
% @inproceedings{csam2024, ... }
% @inproceedings{csam2validation2024, ... }
% @article{hallucinationsurvey, ... }
% @inproceedings{factscore2023, ... }
% @inproceedings{halueval2023, ... }
% @inproceedings{ragoriginal, ... }
% @article{ragsurvey2024, ... }
% @article{frameworkcomparison2024, ... }
% @article{llmsurvey2023, ... }
% @book{dsstheory, ... }
% @article{vectordbsurvey, ... }

\end{document}
